\documentclass[10pt,a4paper]{amsart}
\usepackage[margin=19mm]{geometry}
\usepackage{amsmath,amssymb,mathtools}
\usepackage[scr=rsfs,cal=euler]{mathalfa}
\usepackage{xfrac}
\usepackage[unicode,hidelinks,bookmarks=false]{hyperref}
\newcommand{\ie}{i.e.\ }
\newcommand{\cf}{cf.\ }
\newcommand{\resp}{respectively\ }
\newcommand{\Subsection}{\subsection}
\providecommand{\nobreakdash}{\nobreak\hbox{-}}
\setlength{\emergencystretch}{2em}
\multlinegap0pt
\usepackage{tikz}
\usetikzlibrary{cd}
\usepackage{mathtools}
\usepackage{amssymb}
\usepackage{xcolor}
\usepackage{enumerate}
\newtheorem{theorem}{Theorem}
\newtheorem{lemma}{Lemma}
\newcommand{\NN}{{\mathbb N}}
\newcommand{\ZZ}{{\mathbb Z}}
\def\a{{\bold a}}
\def\b{{\bold b}}
\def\e{{\bold e}}
\def\kk{{\Bbbk}}
\def\mm{{\mathfrak m}}
\def\x{{\bold x}}
\title{Relative Hochster--Takayama formula and Cohen--Macaulay monomial ideal quotients}
\author{T\`ai Huy H\`a, Nguyen Cong Minh}
\date{Source: arXiv:2601.20233v1 (CC BY 4.0)}
\begin{document}
\maketitle

\noindent\emph{Source and licence.} Relative Hochster--Takayama formula and Cohen--Macaulay monomial ideal quotients. Version arXiv:2601.20233v1 (CC BY 4.0), distributed by arXiv under the Creative Commons Attribution 4.0 International licence, http://creativecommons.org/licenses/by/4.0/, as recorded in the arXiv licence metadata for this exact version and shown on the abstract page https://arxiv.org/abs/2601.20233v1. Full licence notice: \texttt{licenses/arxiv-2601.20233-CC-BY-4.0.txt}.\par\smallskip

\noindent\textit{Editorial scope.} Counted unit: the article's theorem MainLem (source0.tex lines 504--541). The article's complete original proof of it is delivered with it (source0.tex lines 543--629, one \texttt{proof} environment of 87 lines). No mathematics is written, repaired or re-derived: the statement and the proof are the article's own lines, in the article's own notation, and no retained line is substituted or re-flowed.  Retained uncounted as the source-local context the proof needs: lines 314--318; lines 422--437.  Nothing was omitted from any retained span, so the package carries no omission record.  No retained line contains a citation, so the wrapper carries no bibliography and no bibliography entry.  The retained lines contain the article's own diagram code, which the wrapper typesets with the standard tikz packages; no image file is delivered and the package declares no asset dependency.  Editorial formatting only, in addition: an AMS \texttt{amsart} preamble that restates the article's own declarations and macros used by the retained lines, namely the environments \texttt{theorem}, \texttt{lemma} on separate counters, together with the standard packages those lines need.  The retained environments print in this document as Theorem 1, Lemma 1 in order of appearance, whereas the article prints the same items with the numbering of its own full document.  The complete native source of the article as distributed in the arXiv e-print for this exact version is bundled unchanged as \texttt{sources/193512/source0.tex}.
\begin{theorem}[Relative Hochster--Takayama]\label{HochsterTakayama}
	Let $J\subsetneq I$ be monomial ideals in $S$. For each $i\in\ZZ$ and $\a\in\ZZ^n$, there is an isomorphism of $\kk$-vector spaces
	$$H_{\mm}^{i}(I/J)_{\a}\cong \widetilde{H}^{i-|G_\a|-1}((\Delta_\a(J),\Delta_\a(I)); \kk),$$
	and $H_{\mm}^{i}(I/J)_{\a}=(0)$ if $G_\a \not\in \Delta(\sqrt{J})$.
\end{theorem}

\begin{lemma}\label{lem:gradedpiece-dim1}
	Let $J\subseteq I$ be monomial ideals in $S$, and set $M:=I/J$.
	Fix $F\subseteq[n]$ and $\a = (a_1, \dots, a_n) \in\ZZ^n$. Write $\a^-:=\sum_{i\in G_\a}(-a_i)\e_i\in\NN^n \text{ and } \a^+:=\a+\a^-\in\NN^n.$
	The following statements hold:
	\begin{enumerate}
		\item[(1)] If $F\not\supseteq G_\a$, then $(M_{x_F})_\a=0$.
		\item[(2)] If $F\supseteq G_\a$, then $(M_{x_F})_\a$ is either $0$ or a $1$-dimensional
		$k$-vector space. More precisely,
		\[
		(M_{x_F})_\a\neq 0
		\quad\Longleftrightarrow\quad
		x^{\a^+}\in I S_{x_F}\setminus J S_{x_F},
		\]
		and in this case $(M_{x_F})_\a$ is generated by the class of $\frac{x^{\a^+}}{x^{\a^-}}\ \in\ M_{x_F}.$
	\end{enumerate}
\end{lemma}

\begin{theorem}[Functoriality]\label{MainLem}
	Let $J\subseteq I$ be monomial ideals in $S$ and let $M = I/J$.
	Fix $i\in\mathbb Z_{\ge 0}$, and let $\a\in\mathbb Z^n$ and $\b\in\mathbb N^n$ be such that $G_{\a+\b} = G_\a$.
	Then multiplication by $\x^\b$ induces a commutative diagram
	\[
	\begin{tikzcd}
		H^i_{\mathfrak m}(M)_\a
		\arrow[r,"{\cdot \x^\b}"]
		\arrow[d,"\cong"']
		&
		H^i_{\mathfrak m}(M)_{\a+\b}
		\arrow[d,"\cong"]
		\\
		\widetilde H^{\,i-|G_\a|-1}\!
		\big(
		\Delta_\a(J),\,\Delta_\a(I)
		\big)
		\arrow[r]
		&
		\widetilde H^{\,i-|G_\a|-1}\!
		\big(
		\Delta_{\a+\b}(J),\,\Delta_{\a+\b}(I)
		\big),
	\end{tikzcd}
	\]
	where the vertical isomorphisms are given by the relative Hochster--Takayama
	formula, and the bottom horizontal map is induced by the natural inclusions
	\[
	\Delta_{\a+\b}(J)\ \subseteq\ \Delta_\a(J)
	\text{ and }
	\Delta_{\a+\b}(I)\ \subseteq\ \Delta_\a(I).
	\]
	In particular, under the hypothesis $G_{\a+\b}=G_\a$, the multiplication map
	\[
	\cdot \x^\b:\ H^i_{\mathfrak m}(M)_\a \longrightarrow H^i_{\mathfrak m}(M)_{\a+\b}
	\]
	is zero if and only if the induced map on relative reduced cohomology is zero.
\end{theorem}

\begin{proof}
	Let $\b\in\NN^n$ and $\a\in\ZZ^n$ satisfy $G_{\a+\b}=G_\a$. Set $G:=G_\a=G_{\a+\b}$.
	Multiplication by $x^\b$ defines a homogeneous $S$-linear map $M\to M$ of degree $\b$, so a morphism of
	$\ZZ^n$-graded \v{C}ech complexes
	\[
	\mu_\b^\bullet:\ \Check{\bf C}^\bullet\otimes_S M \longrightarrow \Check{\bf C}^\bullet\otimes_S M,
	\qquad
	u\longmapsto x^\b u,
	\]
	which commutes with the \v{C}ech differential because all localization maps are $S$-linear. Taking the $\a$-graded
	strand gives a cochain map
	\[
	(\mu_\b^\bullet)_\a:\ (\Check{\bf C}^\bullet\otimes_S M)_\a \longrightarrow (\Check{\bf C}^\bullet\otimes_S M)_{\a+\b},
	\]
	whose induced map on cohomology is exactly $\cdot x^\b:\ H^i_\mm(M)_\a\to H^i_\mm(M)_{\a+\b}$.
	
	For any monomial ideal $K \subseteq S$, we observe that $\Delta_{\a+\b}(K)\subseteq \Delta_\a(K)$.
	Indeed, let $E\subseteq[n]\setminus G$ and set $F:=G\cup E$. It follow from the definition of degree complexes that
	\[
	E\in\Delta_{\a}(K)\quad\Longleftrightarrow\quad x^{\a^+}\notin K S_{x_F}\qquad(\a\in\ZZ^n).
	\]
	If $E\in\Delta_{\a+\b}(K)$, then $x^{(\a+\b)^+}=x^{\a^++\b}\notin K S_{x_F}$. If $x^{\a^+}\in K S_{x_F}$, then multiplying
	by $x^\b\in S\subseteq S_{x_F}$ gives $x^{\a^++\b}\in K S_{x_F}$, a contradiction. Hence, $x^{\a^+}\notin K S_{x_F}$ and
	$E\in\Delta_\a(K)$. This proves the inclusions for $K=I,J$, hence an inclusion of pairs
	\[
	(\Delta_{\a+\b}(J),\Delta_{\a+\b}(I))\hookrightarrow (\Delta_\a(J),\Delta_\a(I)).
	\]
	Let $\iota^\ast$ denote the induced restriction map on reduced relative cohomology.
	
	Recall from the proof of Theorem~\ref{HochsterTakayama} that the isomorphism
	\[
	H^i_\mm(M)_\a\ \cong\ \widetilde H^{\,i-|G_\a|-1}\big(\Delta_\a(J),\Delta_\a(I);\kk\big)
	\]
	is induced by an identification of cochain complexes
	\[
	\Phi_\a:\ (\Check{\bf C}^\bullet\otimes_S M)_\a[|G_\a|+1]\ \xrightarrow{\ \cong\ }\
	\widetilde C^\bullet\big(\Delta_\a(J),\Delta_\a(I);\kk\big),
	\]
	whose natural basis elements are described as follows.
	Fix $E\subseteq[n]\setminus G_\a$ and set $F:=G_\a\cup E$. By Lemma~\ref{lem:gradedpiece-dim1},
	$(M_{x_F})_\a$ is $0$ or $1$-dimensional, and if it is nonzero then it is generated by the class
	\[
	u_E(\a):=\left[\frac{x^{\a^+}}{x^{\a^-}}\right]\in (M_{x_F})_\a.
	\]
	Moreover, $(M_{x_F})_\a\neq 0$ holds if and only if $E\in\Delta_\a(J)\setminus\Delta_\a(I)$.
	Under $\Phi_\a$, the generator $u_E(\a)$ corresponds to the standard basis cochain $\varphi_E$ of the relative
	cochain complex of $(\Delta_\a(J),\Delta_\a(I))$.
	
	Now, fix $E\subseteq[n]\setminus G$ and $F:=G\cup E$. Assume $(M_{x_F})_\a\neq 0$, so $E\in\Delta_\a(J)\setminus\Delta_\a(I)$
	and $u_E(\a)$ is defined. Because $G_{\a+\b}=G_\a$, we have $(\a+\b)^-=\a^-$ and $(\a+\b)^+=\a^++\b$.
	Thus in $M_{x_F}$,
	\[
	x^\b\cdot u_E(\a)
	=
	x^\b\cdot \left[\frac{x^{\a^+}}{x^{\a^-}}\right]
	=
	\left[\frac{x^{\a^++\b}}{x^{\a^-}}\right]
	=
	\left[\frac{x^{(\a+\b)^+}}{x^{(\a+\b)^-}}\right]
	=
	u_E(\a+\b).
	\]
	If $E\in\Delta_{\a+\b}(J)\setminus\Delta_{\a+\b}(I)$, then $(M_{x_F})_{\a+\b}\neq 0$ and $u_E(\a+\b)$ is the generator,
	so $(\mu_\b^\bullet)_\a$ sends $u_E(\a)$ to the generator corresponding to the same face $E$.
	If $E\notin\Delta_{\a+\b}(J)\setminus\Delta_{\a+\b}(I)$, then $(M_{x_F})_{\a+\b}=0$ and hence $(\mu_\b^\bullet)_\a$ sends
	$u_E(\a)$ to $0$ in degree $\a+\b$.
	
	On the simplicial side, the restriction map $\iota^\ast$ sends the basis cochain $\varphi_E$ to $\varphi_E$ if
	$E$ is a face of the smaller relative complex, and to $0$ otherwise. We conclude, therefore, that the diagram of cochain complexes
	\[
	\begin{tikzcd}
		(\Check{\bf C}^\bullet\otimes_S M)_\a[|G|+1]
		\arrow[r,"(\mu_\b^\bullet)_\a"]
		\arrow[d,"\Phi_\a"']
		&
		(\Check{\bf C}^\bullet\otimes_S M)_{\a+\b}[|G|+1]
		\arrow[d,"\Phi_{\a+\b}"]
		\\
		\widetilde C^\bullet(\Delta_\a(J),\Delta_\a(I);\kk)
		\arrow[r,"\iota^\ast"]
		&
		\widetilde C^\bullet(\Delta_{\a+\b}(J),\Delta_{\a+\b}(I);\kk)
	\end{tikzcd}
	\]
	commutes on basis elements, and so commutes.
	Passing to cohomology gives the desired commutative diagram.
\end{proof}

\end{document}
